\editionPage[\pdfbookmark[0]{Deux bilans annuels}{bm-bilans-guide-parent}]{bilans-guide}{Deux bilans annuels}{Sujets originaux · Réviser en reliant les notions}{Durée conseillée : 2 heures · Barème d’entraînement sur 20}
\begin{blocEdition}{bilans-guide}

Ces évaluations sont des créations du livre pour la Cinquième. Elles n’ont
pas de statut d’examen officiel. Le temps et le barème ci-dessous sont des
repères pour une séance autonome ; ils peuvent être adaptés en classe.
\editionEtape{Automatismes : environ 20 minutes}
Douze questions courtes, conseillées sans calculatrice : 6 points, soit
0,5 point par réponse. Écrire les étapes utiles, même lorsqu’elles sont brèves.
\editionEtape{Problèmes : environ 1 h 40}
Quatre problèmes indépendants, calculatrice autorisée : 14 points.
Chaque problème vaut 3,5 points : trois questions à 1 point et 0,5 point
pour la clarté de la rédaction, les unités et les contrôles.
\editionEtape{Travailler après l’évaluation}
Comparer les démarches, relever les erreurs par notion, puis refaire
les questions utiles sans correction. Un résultat faux peut accompagner
une démarche partiellement juste : identifier précisément l’étape à reprendre.
\editionEtape{Préparer une autre séance}
Revenir au cours et à deux exercices du chapitre concerné. Le second
bilan peut être fait plus tard pour vérifier que la méthode est devenue disponible.

\end{blocEdition}
\editionPage{bilan-A-auto}{Bilan A}{Relier calculs, données et géométrie}{Partie 1 · 20 min · Sans calculatrice · 6 points}
\begin{blocEdition}{bilan-A-auto}
\begin{questions}
\item Calculer $7+4\times3$.
\item Calculer $8{,}4:0{,}7$.
\item Donner quotient et reste de 74 par 9.
\item Comparer $-3{,}4$ et $-3{,}04$.
\item Calculer $-7-(-12)$.
\item Calculer $2/3+1/4$.
\item Calculer $11^2$.
\item Développer $4(x+2)$.
\item Résoudre $5x=18$.
\item Calculer 25 \% de 60.
\item Deux angles d’un triangle valent 55 et 65 degrés. Trouver le troisième.
\item Quelle est la probabilité d’obtenir un nombre impair sur un dé équilibré à six faces ?\end{questions}
\end{blocEdition}
\editionPage{bilan-A-ex-1}{Bilan A · Problème 1}{Une sortie et son budget}{Partie 2 · 1 h 40 pour les quatre problèmes · 3,5 points}
\begin{blocEdition}{bilan-A-ex-1}
\begin{questions}
\item 28 élèves paient chacun 8,50 euros. Le transport coûte 98 euros et les entrées 4,50 euros par élève. Calculer la dépense totale.
\item Quelle somme reste-t-il ? Quel pourcentage des sommes reçues cela représente-t-il ?
\item Le reste est partagé également entre sept ateliers. Combien reçoit chacun ?
\end{questions}\par\vspace{5mm}{\otLight\fontsize{8}{10}\selectfont\color{otMuted}Recherche et rédaction}\par\par\vspace{7mm}{\color{otRule!65}\hrule height .25pt}\par\vspace{7mm}{\color{otRule!65}\hrule height .25pt}\par\vspace{7mm}{\color{otRule!65}\hrule height .25pt}\par\vspace{7mm}{\color{otRule!65}\hrule height .25pt}\par\vspace{7mm}{\color{otRule!65}\hrule height .25pt}\par\vspace{7mm}{\color{otRule!65}\hrule height .25pt}
\end{blocEdition}
\editionPage{bilan-A-ex-2}{Bilan A · Problème 2}{Un triangle et sa médiane}{Partie 2 · 1 h 40 pour les quatre problèmes · 3,5 points}
\begin{blocEdition}{bilan-A-ex-2}
\begin{questions}
\item ABC a $AB=10$ cm et une hauteur issue de C de 6 cm. Calculer son aire.
\item M est le milieu de [AB]. Calculer l’aire de ACM et justifier.
\item N est le milieu de [AM]. Calculer l’aire de ACN, puis celle de NCM.
\end{questions}\label{fig-bilan-A-ex-2}\figureGeo{bilan-triangle}
\par\vspace{5mm}{\otLight\fontsize{8}{10}\selectfont\color{otMuted}Recherche et rédaction}\par\par\vspace{7mm}{\color{otRule!65}\hrule height .25pt}\par\vspace{7mm}{\color{otRule!65}\hrule height .25pt}\par\vspace{7mm}{\color{otRule!65}\hrule height .25pt}\par\vspace{7mm}{\color{otRule!65}\hrule height .25pt}\par\vspace{7mm}{\color{otRule!65}\hrule height .25pt}\par\vspace{7mm}{\color{otRule!65}\hrule height .25pt}
\end{blocEdition}
\editionPage{bilan-A-ex-3}{Bilan A · Problème 3}{Une enquête et un tirage}{Partie 2 · 1 h 40 pour les quatre problèmes · 3,5 points}
\begin{blocEdition}{bilan-A-ex-3}
\begin{questions}
\item Une enquête auprès de 20 élèves donne 8 vélos, 6 bus et 6 déplacements à pied. Donner les trois fréquences.
\item Donner les angles d’un diagramme circulaire représentant cette enquête.
\item On tire au hasard une fiche parmi les 20 fiches individuelles, toutes de même taille. Donner la probabilité que la fiche corresponde à un élève venu en bus, puis expliquer le modèle.
\end{questions}\label{fig-bilan-A-ex-3}\figureGeo{bilan-enquete}
\par\vspace{5mm}{\otLight\fontsize{8}{10}\selectfont\color{otMuted}Recherche et rédaction}\par\par\vspace{7mm}{\color{otRule!65}\hrule height .25pt}\par\vspace{7mm}{\color{otRule!65}\hrule height .25pt}\par\vspace{7mm}{\color{otRule!65}\hrule height .25pt}\par\vspace{7mm}{\color{otRule!65}\hrule height .25pt}\par\vspace{7mm}{\color{otRule!65}\hrule height .25pt}\par\vspace{7mm}{\color{otRule!65}\hrule height .25pt}
\end{blocEdition}
\editionPage{bilan-A-ex-4}{Bilan A · Problème 4}{Un bac et des récipients}{Partie 2 · 1 h 40 pour les quatre problèmes · 3,5 points}
\begin{blocEdition}{bilan-A-ex-4}
\begin{questions}
\item Un bac intérieur mesure 40 cm sur 25 cm sur 30 cm. Calculer son volume en cm$^3$ puis en litres.
\item On remplit seulement les $3/5$ du bac. Quel volume d’eau cela représente-t-il ?
\item On dispose de bouteilles de 1,5 L. Combien de bouteilles pleines faut-il pour ce remplissage ?
\end{questions}\label{fig-bilan-A-ex-4}\figureGeo{bilan-bac}
\par\vspace{5mm}{\otLight\fontsize{8}{10}\selectfont\color{otMuted}Recherche et rédaction}\par\par\vspace{7mm}{\color{otRule!65}\hrule height .25pt}\par\vspace{7mm}{\color{otRule!65}\hrule height .25pt}\par\vspace{7mm}{\color{otRule!65}\hrule height .25pt}\par\vspace{7mm}{\color{otRule!65}\hrule height .25pt}\par\vspace{7mm}{\color{otRule!65}\hrule height .25pt}\par\vspace{7mm}{\color{otRule!65}\hrule height .25pt}
\end{blocEdition}
\editionPage{bilan-B-auto}{Bilan B}{Chercher, modéliser et expliquer}{Partie 1 · 20 min · Sans calculatrice · 6 points}
\begin{blocEdition}{bilan-B-auto}
\begin{questions}
\item Calculer $24:3\times2$.
\item 315 est-il divisible par 9 ?
\item Donner l’opposé de $-2{,}5$ et sa distance à zéro.
\item Calculer $-3+8-6$.
\item Comparer $3/4$ et $5/8$.
\item Calculer $5/6-1/4$.
\item Calculer $3^3+2^2$.
\item Réduire $2x+5+3x+1$.
\item Résoudre $x+7=3$.
\item Convertir 4500 cm$^3$ en litres.
\item Calculer la moyenne de 6, 9 et 12.
\item Un demi-tour conserve-t-il les aires ?\end{questions}
\end{blocEdition}
\editionPage{bilan-B-ex-1}{Bilan B · Problème 1}{Deux tarifs de location}{Partie 2 · 1 h 40 pour les quatre problèmes · 3,5 points}
\begin{blocEdition}{bilan-B-ex-1}
\begin{questions}
\item A coûte 4 euros par heure. B coûte 6 euros fixes puis 2 euros par heure. Construire un tableau pour 0, 1, 2, 3 et 4 h.
\item Lire quand les deux tarifs sont égaux. Comparer pour 2 h puis 4 h.
\item Quel tarif est proportionnel à la durée ? Justifier par le tableau et le graphique.
\end{questions}\label{fig-bilan-B-ex-1}\figureGeo{bilan-tarifs}
\par\vspace{5mm}{\otLight\fontsize{8}{10}\selectfont\color{otMuted}Recherche et rédaction}\par\par\vspace{7mm}{\color{otRule!65}\hrule height .25pt}\par\vspace{7mm}{\color{otRule!65}\hrule height .25pt}\par\vspace{7mm}{\color{otRule!65}\hrule height .25pt}\par\vspace{7mm}{\color{otRule!65}\hrule height .25pt}\par\vspace{7mm}{\color{otRule!65}\hrule height .25pt}\par\vspace{7mm}{\color{otRule!65}\hrule height .25pt}
\end{blocEdition}
\editionPage{bilan-B-ex-2}{Bilan B · Problème 2}{Un parallélogramme et ses diagonales}{Partie 2 · 1 h 40 pour les quatre problèmes · 3,5 points}
\begin{blocEdition}{bilan-B-ex-2}
\begin{questions}
\item Dans ABCD, O est le milieu de [AC] et de [BD]. Quelle nature peut-on prouver ?
\item On sait en plus que $AC=BD$. Que peut-on préciser ? Et si les diagonales sont aussi perpendiculaires ?
\item Pour $AB=8$ cm et $BC=5$ cm dans le cas du rectangle, calculer aire et périmètre.
\end{questions}\label{fig-bilan-B-ex-2}\figureGeo{bilan-rectangle}
\par\vspace{5mm}{\otLight\fontsize{8}{10}\selectfont\color{otMuted}Recherche et rédaction}\par\par\vspace{7mm}{\color{otRule!65}\hrule height .25pt}\par\vspace{7mm}{\color{otRule!65}\hrule height .25pt}\par\vspace{7mm}{\color{otRule!65}\hrule height .25pt}\par\vspace{7mm}{\color{otRule!65}\hrule height .25pt}\par\vspace{7mm}{\color{otRule!65}\hrule height .25pt}\par\vspace{7mm}{\color{otRule!65}\hrule height .25pt}
\end{blocEdition}
\editionPage{bilan-B-ex-3}{Bilan B · Problème 3}{Un cylindre à remplir}{Partie 2 · 1 h 40 pour les quatre problèmes · 3,5 points}
\begin{blocEdition}{bilan-B-ex-3}
\begin{questions}
\item Un cylindre a un rayon intérieur de 3 cm et une hauteur de 12 cm. Calculer son volume exact.
\item Arrondir sa capacité au mL. Expliquer la conversion.
\item Un second cylindre a le même rayon mais une hauteur de 8 cm. Donner le rapport de sa capacité à celle du premier et l’exprimer en pourcentage, au dixième près.
\end{questions}\label{fig-bilan-B-ex-3}\figureGeo{bilan-cylindre}
\par\vspace{5mm}{\otLight\fontsize{8}{10}\selectfont\color{otMuted}Recherche et rédaction}\par\par\vspace{7mm}{\color{otRule!65}\hrule height .25pt}\par\vspace{7mm}{\color{otRule!65}\hrule height .25pt}\par\vspace{7mm}{\color{otRule!65}\hrule height .25pt}\par\vspace{7mm}{\color{otRule!65}\hrule height .25pt}\par\vspace{7mm}{\color{otRule!65}\hrule height .25pt}\par\vspace{7mm}{\color{otRule!65}\hrule height .25pt}
\end{blocEdition}
\editionPage{bilan-B-ex-4}{Bilan B · Problème 4}{Un script à lire et à réparer}{Partie 2 · 1 h 40 pour les quatre problèmes · 3,5 points}
\begin{blocEdition}{bilan-B-ex-4}
\begin{questions}
\item Le script donné répète trois fois : avancer de 40 pas, tourner à gauche de 90 degrés. Décrire la figure obtenue.
\item Corriger un seul paramètre pour obtenir un carré et donner son périmètre.
\item On veut ensuite un carré de périmètre 240 pas. Quel paramètre modifier ? Pourquoi un virage de 60 degrés ne convient-il pas pour un carré ?
\end{questions}\label{fig-bilan-B-ex-4}\figureGeo{bilan-script}
\par\vspace{5mm}{\otLight\fontsize{8}{10}\selectfont\color{otMuted}Recherche et rédaction}\par\par\vspace{7mm}{\color{otRule!65}\hrule height .25pt}\par\vspace{7mm}{\color{otRule!65}\hrule height .25pt}\par\vspace{7mm}{\color{otRule!65}\hrule height .25pt}\par\vspace{7mm}{\color{otRule!65}\hrule height .25pt}\par\vspace{7mm}{\color{otRule!65}\hrule height .25pt}\par\vspace{7mm}{\color{otRule!65}\hrule height .25pt}
\end{blocEdition}
